% ============================================================================
% Sample gain as a function of the study-only loss correlation
% Draft fragment for PAPER.tex (notation of Sec. 3 and of the appendix
% "Estimating inter-fold covariance in cross-validation"). Written 2026-09-29.
% Needs amsmath. No new macros. Empirical numbers refer to benchmark_regression
% analysis outputs (see the comment block at the end).
% ============================================================================

\subsection{From the fold-error correlation to a study-only statistic}
\label{sec:gain-vs-loss-correlation}

\paragraph{Setting.}
A study set of $N$ observations is split $K$ times by MCCV: each split $a$ draws a test
set $\mathcal{D}^{\text{te}}_a$ of size $n_{\text{te}}$ uniformly without replacement,
independently across splits, and trains $g_a = F_\lambda(\mathcal{D}^{\text{tr}}_a,\xi_a)$
on the remaining $n_{\text{tr}} = N - n_{\text{te}}$ observations. We write
$t = n_{\text{te}}/N$ for the test fraction. For an observation $i$ held out by split $a$,
let
\begin{equation}
  d_{a,i} = \ell\bigl(g_a(x_i), y_i\bigr) - \mathcal{R}^*(g_a)
\end{equation}
be its per-sample loss centred at the risk of the predictor that did not see it, so that
$\mathbb{E}[d_{a,i} \mid g_a] = 0$, and let $\sigma^2_e = \mathbb{E}[d_{a,i}^2]$ be the
per-sample loss variance.

\paragraph{Step 1: the gain depends on the fold-error correlation only.}
Neglecting the evaluation noise of the benchmark estimate (eq.~\eqref{eq:delta_variance_approx}),
the benchmark-adjusted fold error is the test-set average of the centred losses,
\begin{equation}
  \delta_a = \frac{1}{n_{\text{te}}} \sum_{i \in \mathcal{D}^{\text{te}}_a} d_{a,i}.
\end{equation}
By eq.~\eqref{eq:delta_cvvar} and the local $1/m$ scaling of the single-split variance
(eq.~\eqref{eq:sample-gain-def}), the sample gain has the intraclass-correlation form of
eq.~\eqref{eq:sample-gain-icc}:
\begin{equation}
  G^{\text{test}}_K \approx \frac{K}{1 + (K-1)\,\rho_\delta},
  \qquad
  \rho_\delta = \frac{\tau^{\mathrm{te}}}{\sigma^{2,\mathrm{te}}}
              = \operatorname{Corr}(\delta_a, \delta_b), \; a \neq b.
  \label{eq:gain-rho-delta}
\end{equation}
This step is exact up to the $1/m$ scaling; everything the gain depends on is in $\rho_\delta$.

\paragraph{Step 2: decomposing the fold-error covariance.}
The single-split variance is $\sigma^{2,\mathrm{te}} = \sigma^2_e / n_{\text{te}}$: given
$g_a$ the held-out observations are i.i.d.\ and the $d_{a,i}$ are centred, so cross terms
within a split vanish. For two splits $a \neq b$, expanding the product of the two averages
and separating the observations held out by both,
$I_{ab} = \mathcal{D}^{\text{te}}_a \cap \mathcal{D}^{\text{te}}_b$, from the pairs of distinct
observations,
\begin{equation}
  \tau^{\mathrm{te}}
  = \frac{1}{n_{\text{te}}^2}
    \Bigl(
      \underbrace{\mathbb{E}\Bigl[\sum_{i \in I_{ab}} d_{a,i}\, d_{b,i}\Bigr]}_{\text{shared observations}}
      +
      \underbrace{\mathbb{E}\Bigl[\sum_{i \in \mathcal{D}^{\text{te}}_a}\;
                  \sum_{j \in \mathcal{D}^{\text{te}}_b,\, j \neq i} d_{a,i}\, d_{b,j}\Bigr]}_{\text{train/test coupling}}
    \Bigr).
  \label{eq:tau-te-split}
\end{equation}

\emph{Shared observations.} Two independent uniform test sets share on average
$\mathbb{E}|I_{ab}| = n_{\text{te}}^2 / N = t\, n_{\text{te}}$ observations. On each of them
the two split predictors' losses covary by $\rho_e\,\sigma^2_e$, where
\begin{equation}
  \rho_e = \frac{\mathbb{E}[d_{a,i}\, d_{b,i} \mid i \in I_{ab}]}{\sigma^2_e}
\end{equation}
is the correlation of the per-sample losses of two split predictors on an observation that
neither was trained on. The first term of~\eqref{eq:tau-te-split} is therefore
$t\, n_{\text{te}}\, \rho_e\, \sigma^2_e / n_{\text{te}}^2$.

\emph{Train/test coupling.} For $i \neq j$, $d_{a,i}$ and $d_{b,j}$ are not independent: $i$
may belong to the training set of $g_b$, $j$ to that of $g_a$, and $g_a$, $g_b$ share most of
their training data. This is intrinsic to cross-validation and is part of $\tau^{\mathrm{te}}$
(Sec.~\ref{sec:sample_gain}). Writing $\gamma = \mathbb{E}[d_{a,i}\, d_{b,j}]$ for the average
covariance of such a pair, the second term is approximately $\gamma$ (about $n_{\text{te}}^2$ pairs).

Dividing by $\sigma^{2,\mathrm{te}} = \sigma^2_e / n_{\text{te}}$,
\begin{equation}
  \boxed{\;
  \rho_\delta = t\,\rho_e + \kappa,
  \qquad
  \kappa = \frac{n_{\text{te}}\,\gamma}{\sigma^2_e}
  \;}
  \label{eq:rho-delta-decomp}
\end{equation}
The loss correlation $\rho_e$ measures how similarly two split predictors err on the same
observation; the factor $t$ is the share of a test set made of observations that another split
also tests; $\kappa$ is the coupling contribution, of either sign.

\paragraph{Step 3: the reference curve and the study-only statistic.}
Substituting~\eqref{eq:rho-delta-decomp} into~\eqref{eq:gain-rho-delta},
\begin{equation}
  G^{\text{test}}_K \approx \frac{K}{1 + (K-1)\,(t\,\rho_e + \kappa)}.
  \label{eq:gain-full}
\end{equation}
With $\kappa = 0$ this is the reference curve $K/(1+(K-1)x)$ with $x = t\rho_e$. Two limits
read directly: for $\rho_e \to 0$ (split predictors err independently) $G^{\text{test}}_K \to K$;
for $\rho_e \to 1$ (all split predictors err alike) $G^{\text{test}}_K \to K/(1+(K-1)t)$, whose
$K\to\infty$ limit $1/t = N/n_{\text{te}}$ says that CV then averages over the whole study set
and no more.

The study-only estimate of $x$ uses the loss-correlation factor of the redundancy score of
Sec.~\ref{sec:study-only-redundancy-statistics}, computed on the per-sample loss of the metric
whose gain is targeted, after the first $k$ splits of a single run:
\begin{equation}
  \widehat{x}_{k} = t\, \widehat\rho^{\mathrm{study}}_{e,k},
  \qquad
  \widehat\rho^{\mathrm{study}}_{e,k}
  = \frac{\widehat C^{\mathrm{study}}_{e,k}}{\widehat V^{\mathrm{study}}_{e,k}}
  \quad \text{(eq.~\eqref{eq:study-loss-correlation})}.
  \label{eq:x-hat}
\end{equation}
The empirical overlap $\bar m^{\mathrm{study}}_k / n_{\text{te}}$ may replace $t$ (same expectation).

\paragraph{Relation to the redundancy score.}
Since $\mathbb{E}\,\bar m^{\mathrm{study}}_k = t\, n_{\text{te}}$,
$\widehat\omega^{\mathrm{study}}_k = \widehat C^{\mathrm{study}}_{g,k}\,
\widehat\rho^{\mathrm{study}}_{e,k}\, \bar m^{\mathrm{study}}_k
\approx \widehat C^{\mathrm{study}}_{g,k}\, n_{\text{te}}\; \widehat{x}_k$.
The score thus contains $\widehat{x}_k$, multiplied by the prediction covariance (which carries
the squared units of the target and the signal strength) and by the test-set size (which grows
with the study). Both factors vary across settings independently of the gain, which is why
$\widehat{x}_k$, dimensionless and bounded by $t$, is the quantity on which an absolute
early-stopping threshold can be placed.

\paragraph{Assumptions and scope.}
(a) The benchmark noise is negligible and the single-split variance scales as $1/m$
(Sec.~\ref{sec:sample_gain}). (b) $\widehat\rho^{\mathrm{study}}_{e,k}$, computed on
$I_{ab}$ with sample centring, estimates $\rho_e$; its bias is $O(1/|I_{ab}|)$, which requires
held-out folds of a few observations (Sec.~\ref{sec:study-only-redundancy-statistics}).
(c) The overlap $\mathbb{E}|I_{ab}| = t\,n_{\text{te}}$ holds for MCCV; for repeated $K$-fold,
splits of the same repeat have disjoint test sets and the factor changes accordingly.
(d) The reference curve omits $\kappa$.

\paragraph{Empirical checks (simulated data).}
(1) The study-only $\widehat\rho^{\mathrm{study}}_{e,k}$ matches the population $\rho_e$ of the
same split predictors, measured on the benchmarking set, within $1\%$ across all simulated
designs. (2) Varying the test fraction ($t = 0.1, 0.2, 0.3, 0.5$) with everything else fixed,
the fitted factor $\rho_\delta / \rho_e$ is $0.10, 0.21, 0.32, 0.48$ (classification) and
$0.14, 0.20, 0.36, 0.46$ (regression), with intercepts within $\pm 0.01$. (3) Regression
configurations and the \texttt{PCam} networks lie on the reference curve (\texttt{PCam}:
median $\log(G^{\text{test}}_{20}/\text{curve})$ between $-0.02$ and $+0.04$); simulated
classifiers fall $10$--$50\%$ below it at training sizes $100$--$300$, i.e.\ $\kappa > 0$ there,
and the gap closes by $n_{\text{tr}} \approx 800$.

% ---------------------------------------------------------------------------
% Sources for the numbers (benchmark_regression):
%  (1) analysis/GAIN_DECOMPOSITION.md, term E (study vs population), all ShuffleSplit sets.
%  (2) analysis/testsize_factor_check.py, analysis/out_decomp_testsize/factor_by_test_size.csv
%      (job 240569, seeds 7000-7049, K = 20, train sizes 300 / 1000).
%  (3) analysis/out_pcam_consistency/ (PCam, 20 configurations) and
%      analysis/out_omega_sweep/ (classification sweep, job 231717, seeds 6000-6049).
% ---------------------------------------------------------------------------
